\documentclass[10pt]{article}

% Specify the margins and text size.
\setlength{\textwidth}{6.4in}
\setlength{\textheight}{9.5in}
\setlength{\oddsidemargin}{0pt}
\setlength{\evensidemargin}{0pt}
\setlength{\topmargin}{0pt}
\setlength{\hoffset}{.05in}
\setlength{\voffset}{-.8in}

\setlength{\parskip}{5pt}
\setlength{\parindent}{0pt}

% Load some fonts and symbol packages
\usepackage{latexsym}
\usepackage{pifont}       % contains 'star' symbol for counterinsurgency handbook title
\usepackage{yfonts} 
\usepackage{amsmath}
\usepackage{amsfonts}

\usepackage{graphicx}     % actually, this is loaded by pstricks
\usepackage[T1]{fontenc}
\usepackage{ifthen}
\usepackage{pstricks,pst-func,pst-text,pst-node,multido,pst-plot,calc,pst-3dplot}
%\usepackage[all]{xy}
%\usepackage{animate}

% The hyperref package inserts links into the pdf.
\definecolor{MyLinkColor}{rgb}{.1,.2,1}
\definecolor{MyCiteColor}{rgb}{.1,1,.2}
\definecolor{MyURLColor}{rgb}{.4,.4,.4}
\usepackage[backref=true,pagebackref=false,hyperindex,colorlinks=true,
  linkcolor=MyLinkColor,urlcolor=MyURLColor]{hyperref}


% The tweaklist package is something I found on the web.  It provides a simple interface
% for making changes to spacing used in the itemize and enumerate environments.  Comment
% this out if you don't care to use tweaklist.
\usepackage{tweaklist}
\renewcommand{\itemhook}{\setlength{\parskip}{2pt}\setlength{\parsep}%
{1pt}\setlength{\topsep}{0pt}\setlength{\itemsep}{0pt}}

\newcommand{\N}{\ensuremath\mathbb{N}}
\newcommand{\Q}{\ensuremath\mathbb{Q}}
\newcommand{\R}{\ensuremath\mathbb{R}}
\newcommand{\C}{\ensuremath\mathbb{C}}
\newcommand{\z}{\hphantom{0}}

\begin{document}
\pagestyle{empty}

\href{http://www.su.edu}{\includegraphics[height=1.75cm]{sulogo.eps}}
\vspace{-1.69cm}

{\small\hfill
\begin{tabular}{cl}
 & Math 207\\ &  Introduction to Statistics\\
\end{tabular}
}

\setlength{\baselineskip}{1.05\baselineskip}
\bigskip\bigskip

1. For the following list, convert the numbers to standard units:
1, 3, 4, 4, 5, 7.
\vspace{1.25in}

2. Make a sketch then find the area under the normal curve ---

\begin{center}
\begin{tabular}{lcl}
a) to the right of 1.0  &\hspace{2in} & b) to the left of -1.0\\[1in]
c) to the left of 2.0   &             & d) between -1 and 2
\end{tabular}
\end{center}
\vspace{.75in}

\href{http://www.su.edu}{\includegraphics[height=1.75cm]{sulogo.eps}}
\vspace{-1.69cm}


{\small\hfill
\begin{tabular}{cl}
 & Math 207\\ &  Introduction to Statistics\\
\end{tabular}
}
\bigskip\bigskip

1. For the following list, convert the numbers to standard units:
1, 3, 4, 4, 5, 7.
\vspace{1.25in}

2. Make a sketch then find the area under the normal curve ---

\begin{center}
\begin{tabular}{lcl}
a) to the right of 1.0  &\hspace{2in} & b) to the left of -1.0\\[1in]
c) to the left of 2.0   &             & d) between -1 and 2
\end{tabular}
\end{center}
\eject

3. For the list of numbers below, calculate the 25th percentile, median and 75th percentile.
3, 8, 6, 14, 0, -4, 2, 12, -7, -1, -10.
\vspace{1.25in}

4. True or false:  The smallest the sd can be is 0 but there is not largest possible value.

\vspace{1in}

5. Given that a sample is approximately normal with a mean of 25 and SD of 2, the approximate
percentage of data values that are expected to fall betwen 23 and 27 is 
\underline{\hspace{1in}}.
The approximate
percentage of data values that are expected to fall betwen 19 and 31 is 
\underline{\hspace{1in}}.

\vfill
3. For the list of numbers below, calculate the 25th percentile, median and 75th percentile.
3, 8, 6, 14, 0, -4, 2, 12, -7, -1, -10.
\vspace{1.25in}


4. True or false:  The smallest the sd can be is 0 but there is not largest possible value.
\vspace{1in}

5. Given that a sample is approximately normal with a mean of 25 and SD of 2, the approximate
percentage of data values that are expected to fall betwen 23 and 27 is 
\underline{\hspace{1in}}.
The approximate
percentage of data values that are expected to fall betwen 19 and 31 is 
\underline{\hspace{1in}}.
\vspace{1in}
\end{document}
